\chapter{Developer Documentation}

\section{Introduction}

This chapter provides comprehensive technical documentation for developers who wish to understand, maintain, or extend the NutriFit application. The documentation covers the internal system architecture, database design, implementation details, security mechanisms, and testing methodologies employed in the development of this application. The information presented here assumes familiarity with web development concepts, JavaScript programming, and relational database principles.

The NutriFit application is built using modern web technologies following industry best practices for code organization, security, and performance. The codebase is structured to promote maintainability, scalability, and ease of understanding. This chapter serves as a guide for developers to navigate the codebase effectively and understand the design decisions that shaped the application's architecture.

\section{System Architecture}

\subsection{Overview}

The NutriFit application follows a three-tier architecture pattern that separates concerns into distinct layers, each with specific responsibilities. This architectural approach promotes modularity, maintainability, and scalability by ensuring that changes in one layer have minimal impact on other layers.

The presentation layer, implemented using React and Next.js, handles all user interface rendering and client-side interactions. This layer is responsible for displaying data to users, capturing user inputs, and providing immediate feedback through visual elements. The presentation layer communicates with the application layer through HTTP requests to API endpoints.

The application layer, implemented using Next.js API routes, contains the business logic of the application. This layer processes requests from the presentation layer, performs data validation, executes business rules, and coordinates interactions with the data layer. The application layer ensures that all operations adhere to the application's business requirements and security policies.

The data layer, implemented using Prisma ORM with SQLite database, manages data persistence and retrieval. This layer provides an abstraction over the database, allowing the application layer to interact with data using object-oriented paradigms rather than raw SQL queries. The data layer ensures data integrity through constraints, relationships, and transactions.

The architecture diagram presented in Figure~\ref{fig:architecture} in Chapter 2 illustrates how these three layers interact. Client browsers send HTTPS requests to the Next.js frontend, which renders React components and manages client-side state. When data operations are required, the frontend makes API calls to Next.js API routes. These routes process the requests, interact with Prisma ORM to perform database operations, and return responses to the frontend. Prisma ORM translates high-level operations into SQL queries and executes them against the SQLite database.

\section{Technologies Used}

The technology stack for NutriFit was carefully selected based on several criteria including developer productivity, community support, performance characteristics, and long-term maintainability. This section describes each major technology and the rationale for its selection.

\subsection{Next.js Framework}

Next.js version 14.2.33 serves as the primary framework for both frontend and backend development. Next.js is a React-based framework that provides server-side rendering, static site generation, and API routes in a single integrated package. The framework was chosen for several reasons. First, it offers excellent developer experience with features like hot module replacement, automatic code splitting, and built-in TypeScript support. Second, it provides performance benefits through server-side rendering, which improves initial page load times and search engine optimization. Third, the API routes feature allows backend and frontend code to coexist in the same project, simplifying development and deployment.

\subsection{React Library}

React version 18.3.1 is used for building the user interface components. React's component-based architecture promotes code reusability and maintainability. The library's virtual DOM implementation ensures efficient updates to the user interface, providing smooth user experiences even with complex state changes. React's large ecosystem of third-party libraries and extensive community support make it an ideal choice for building modern web applications.

\subsection{TypeScript Language}

TypeScript version 5.4.5 is used throughout the codebase to add static type checking to JavaScript. TypeScript catches many common programming errors at compile time rather than runtime, significantly reducing bugs in production. The type system also serves as documentation, making it easier for developers to understand function signatures and data structures. TypeScript's integration with modern IDEs provides excellent autocomplete and refactoring capabilities, improving developer productivity.

\subsection{Prisma ORM}

Prisma version 6.17.0 serves as the Object-Relational Mapping layer between the application and the database. Prisma provides a type-safe database client that is automatically generated from the database schema. This approach ensures that database queries are validated at compile time, preventing many common database-related errors. Prisma's migration system manages database schema changes in a controlled, version-controlled manner. The Prisma Studio tool provides a graphical interface for viewing and editing database contents during development.

\subsection{SQLite Database}

SQLite version 3.x is used as the database engine for data persistence. SQLite is a serverless, self-contained database that stores all data in a single file. This simplicity makes it ideal for development and small to medium-scale deployments. SQLite provides full SQL support including transactions, foreign keys, and indexes. While SQLite has limitations for high-concurrency scenarios, it is perfectly adequate for the NutriFit application's requirements and can be replaced with PostgreSQL or MySQL if scaling needs arise.

\subsection{Tailwind CSS}

Tailwind CSS version 3.4.3 is used for styling the user interface. Tailwind is a utility-first CSS framework that provides low-level utility classes for building custom designs. Unlike traditional CSS frameworks that provide pre-designed components, Tailwind allows developers to compose designs using utility classes directly in HTML. This approach results in smaller CSS bundles and more maintainable styles. Tailwind's responsive design utilities make it easy to create interfaces that work well on all device sizes.

\subsection{Additional Libraries}

Several additional libraries support specific functionality. The bcryptjs library version 2.4.3 provides password hashing capabilities using the bcrypt algorithm. The jose library version 5.2.3 handles JSON Web Token generation and verification for authentication. The recharts library version 3.2.1 provides charting components for visualizing progress data. These libraries were chosen for their reliability, security, and active maintenance.

\section{Database Design}

\subsection{Overview}

The database schema for NutriFit consists of eleven interconnected tables that store all application data. The schema is designed following database normalization principles to minimize data redundancy and ensure data integrity. The entity-relationship diagram shown in Figure~\ref{fig:er-diagram} illustrates the relationships between different entities in the system.

\begin{figure}[H]
\centering
\includegraphics[width=0.9\textwidth]{er-diagram.png}
\caption{Entity-Relationship Diagram}
\label{fig:er-diagram}
\end{figure}

Figure~\ref{fig:er-diagram} presents the complete entity-relationship structure of the database. The User entity serves as the central entity, with one-to-one relationship to HealthProfile and one-to-many relationships to all other user-specific entities including Goal, DietPlan, WorkoutPlan, MealLog, ExerciseLog, WaterIntake, and Favorite. The FoodItem and ExerciseItem entities exist independently as catalog tables that are referenced by other entities but do not have direct foreign key relationships.

The class diagram illustrated in Figure~\ref{fig:class-diagram} shows the object-oriented representation of these entities along with their attributes and methods. This diagram provides a developer-centric view of the data model, showing not just the data structure but also the operations that can be performed on each entity.

\begin{figure}[H]
\centering
\includegraphics[width=0.95\textwidth]{class-diagram.png}
\caption{Class Diagram}
\label{fig:class-diagram}
\end{figure}

As shown in Figure~\ref{fig:class-diagram}, the class diagram includes entity classes representing database tables, service classes that implement business logic, and value objects that represent complex data structures. The User class contains attributes for user identification and authentication, along with methods for registration, login, and profile management. The HealthProfile class stores health-related data and includes methods for BMI calculation and calorie calculation. Service classes such as PlanGenerator, SafetyChecker, and ProgressTracker encapsulate complex business logic that operates on multiple entities.

\subsection{Data Integrity}

Data integrity is maintained through several mechanisms implemented at the database level. Primary keys ensure that each record in a table can be uniquely identified. The User table uses an auto-incrementing integer as the primary key, ensuring that each user has a unique identifier. All other tables follow the same pattern with auto-incrementing primary keys.

Foreign key constraints enforce referential integrity between related tables. For example, the HealthProfile table has a foreign key constraint on the userId column that references the User table's primary key. This constraint ensures that a health profile cannot exist without a corresponding user record. The constraint is configured with CASCADE delete behavior, meaning that when a user is deleted, their health profile is automatically deleted as well. This cascading delete behavior is applied to all user-related tables, ensuring that user data is completely removed when an account is deleted.

Unique constraints prevent duplicate data where appropriate. The User table has a unique constraint on the email column, ensuring that each email address can only be registered once. The WaterIntake table has a composite unique constraint on userId and date columns, ensuring that only one water intake record exists per user per day. The Favorite table has a composite unique constraint on userId, itemType, and itemId columns, preventing users from favoriting the same item multiple times.

Check constraints and default values ensure data validity. The HealthProfile table specifies default values for gender, activityLevel, and healthGoals columns, ensuring that these fields always have valid values even if not explicitly provided. The WaterIntake table specifies default values of 0 for glasses and 8 for goal, providing sensible defaults for new records.

Indexes improve query performance for frequently accessed data. The MealLog and ExerciseLog tables have composite indexes on userId and date columns, enabling fast retrieval of logs for a specific user within a date range. The Goal table has a composite index on userId and status columns, allowing efficient filtering of active goals. These indexes significantly improve application performance when querying large datasets.

\section{Application Components}

\subsection{Supporting Modules}

The application is organized into several supporting modules that encapsulate specific functionality. These modules promote code reusability and maintainability by separating concerns into logical units.

The authentication module, located in lib/auth.ts, provides functions for user authentication and authorization. The getUserFromToken function extracts and validates JWT tokens from HTTP request headers, returning the user payload if the token is valid or null if invalid. This function is used by all protected API routes to ensure that only authenticated users can access sensitive operations.

The token handler module, located in lib/tokenHandler.ts, manages JWT token generation and verification. The generateToken function creates a new JWT token with a specified payload and seven-day expiration time. The token is signed using a secret key stored in environment variables. The verifyToken function validates a token's signature and expiration, returning the payload if valid. These functions use the jose library, which provides secure and standards-compliant JWT implementation.

The safety check module, located in lib/safetyCheck.ts, implements health-related calculations and safety validations. The calculateBMI function computes Body Mass Index from weight and height values, returning the BMI value along with a category classification. The calculateDailyCalories function implements the Harris-Benedict equation to determine daily calorie needs based on age, weight, height, gender, activity level, and health goals. The checkExerciseSafety function evaluates whether an exercise is safe for a user based on their medical conditions, returning a safety status and warning message if applicable. The checkFoodSafety function determines whether a food item is safe for a user based on their allergen profile.

The plan generator module, located in lib/planGenerator.ts, contains algorithms for generating personalized diet and workout plans. The generateDietPlan function orchestrates the diet plan creation process by calculating calorie needs, distributing calories across meals, filtering safe foods, and selecting appropriate food combinations. The generateMeal helper function creates individual meals by selecting foods from appropriate categories to ensure balanced nutrition. The generateWorkoutPlan function creates exercise routines by analyzing user medical conditions, determining safe intensity levels, and selecting appropriate exercises.

The database module, located in lib/db.ts, exports a configured Prisma client instance that is used throughout the application for database operations. This module ensures that only one Prisma client instance exists, preventing connection pool exhaustion in serverless environments.

\section{API Routes and Functional Flow}

The application implements RESTful API routes using Next.js API route handlers. Each route is defined in a separate file within the app/api directory structure. The routes follow REST conventions, using HTTP methods to indicate operation types: GET for retrieving data, POST for creating data, PUT for updating data, and DELETE for removing data.

\subsection{Authentication Flow}

The authentication flow begins when a user submits registration or login credentials. Figure~\ref{fig:sequence-auth} illustrates the complete authentication sequence.

\begin{figure}[H]
\centering
\includegraphics[width=0.9\textwidth]{sequence-auth.png}
\caption{Authentication Sequence Diagram}
\label{fig:sequence-auth}
\end{figure}

As shown in Figure~\ref{fig:sequence-auth}, the registration process starts when the user enters their credentials in the browser. The browser sends a POST request to the /api/auth/register endpoint with the user's name, email, and password. The API route validates the input data, checking that all required fields are present and properly formatted. The route then queries the database to verify that the email address is not already registered. If the email is available, the route hashes the password using bcrypt with a salt factor of ten. The hashed password and other user information are stored in the database, creating a new user record. The route then generates a JWT token containing the user's ID and email, signs it with the secret key, and returns it to the browser along with the user information. The browser stores the token in local storage and redirects the user to the dashboard.

The login process follows a similar pattern. The user enters their email and password, which are sent to the /api/auth/login endpoint. The route queries the database to find a user with the provided email address. If found, the route uses bcrypt to compare the provided password with the stored password hash. If the passwords match, the route generates a JWT token and returns it to the browser. If the passwords do not match or the user is not found, the route returns an error message.

For subsequent requests to protected endpoints, the browser includes the JWT token in the Authorization header. The API route extracts the token, verifies its signature and expiration using the verifyToken function, and extracts the user ID from the token payload. This user ID is then used to query or modify user-specific data.

\subsection{Diet Plan Generation Flow}

The diet plan generation process involves multiple steps and calculations. Figure~\ref{fig:sequence-diet-plan} illustrates this process in detail.

\begin{figure}[H]
\centering
\includegraphics[width=0.9\textwidth]{sequence-diet-plan.png}
\caption{Diet Plan Generation Sequence Diagram}
\label{fig:sequence-diet-plan}
\end{figure}

Figure~\ref{fig:sequence-diet-plan} shows that the process begins when the user clicks the "Generate Diet Plan" button in the browser. The browser sends a POST request to the /api/diet-plan/generate endpoint with the authentication token. The API route verifies the token and extracts the user ID. The route then queries the database to retrieve the user's health profile, which contains all necessary information for plan generation.

The route passes the health profile to the PlanGenerator service, which first calculates the user's daily calorie needs using the Harris-Benedict equation. The equation calculates Basal Metabolic Rate (BMR) based on weight, height, age, and gender. For males, the formula is: BMR equals 88.362 plus 13.397 times weight in kilograms plus 4.799 times height in centimeters minus 5.677 times age in years. For females, the formula is: BMR equals 447.593 plus 9.247 times weight plus 3.098 times height minus 4.330 times age. The BMR is then multiplied by an activity level factor: 1.2 for sedentary, 1.375 for light activity, 1.55 for moderate activity, 1.725 for active, and 1.9 for very active. The result is the Total Daily Energy Expenditure (TDEE). This value is adjusted based on health goals: reduced by 500 calories for weight loss or increased by 500 calories for weight gain.

The PlanGenerator then distributes the total calories across four meals: 25 percent for breakfast, 35 percent for lunch, 30 percent for dinner, and 10 percent for snacks. The service queries the database to retrieve all food items, then filters out any foods containing allergens specified in the user's profile. For each meal, the service selects foods from appropriate categories to create balanced nutrition. Breakfast typically includes a protein source, a grain, and a fruit. Lunch includes a protein, a carbohydrate, and two vegetables. Dinner includes a protein, a carbohydrate, and a vegetable. Snacks consist of lighter items such as fruits or nuts.

After generating all meals, the service saves the complete diet plan to the database and returns it to the API route, which sends it back to the browser for display to the user.

\section{Security and Validation}

Security is a critical concern in web applications, and NutriFit implements multiple layers of security to protect user data and prevent unauthorized access.

\subsection{Password Security}

User passwords are never stored in plain text. When a user registers or changes their password, the application uses the bcrypt hashing algorithm to create a secure hash of the password. Bcrypt is specifically designed for password hashing and includes several security features. It automatically generates a random salt for each password, preventing rainbow table attacks. It uses a configurable work factor (set to 10 in this application) that makes the hashing process computationally expensive, slowing down brute force attacks. The work factor can be increased over time as computing power increases, maintaining security without changing the algorithm.

During login, the application hashes the provided password using the same algorithm and compares the result with the stored hash. The bcrypt library provides a compare function that performs this comparison in a timing-safe manner, preventing timing attacks that could reveal information about the password.

\subsection{Authentication Token Security}

JWT tokens are used for maintaining user sessions. Each token contains a payload with the user's ID and email, an expiration time set to seven days from creation, and a signature created using a secret key. The secret key is stored in environment variables and never exposed to clients. The signature ensures that tokens cannot be tampered with; any modification to the token payload invalidates the signature, causing verification to fail.

Tokens are transmitted in the Authorization header using the Bearer scheme, following OAuth 2.0 conventions. This approach is more secure than storing tokens in cookies because it is not vulnerable to Cross-Site Request Forgery (CSRF) attacks. However, tokens stored in local storage are vulnerable to Cross-Site Scripting (XSS) attacks, so the application must ensure that all user inputs are properly sanitized to prevent XSS vulnerabilities.

\subsection{Input Validation}

All user inputs are validated both on the client side and server side. Client-side validation provides immediate feedback to users, improving user experience. Server-side validation is essential for security because client-side validation can be bypassed by malicious users.

The application validates data types, ensuring that numeric fields contain only numbers, email fields contain valid email addresses, and date fields contain valid dates. Range validation ensures that values fall within acceptable ranges; for example, age must be between 1 and 120, weight must be between 20 and 500 kilograms, and height must be between 50 and 300 centimeters.

Required field validation ensures that mandatory fields are not empty. The application checks for the presence of required fields before processing requests and returns clear error messages when fields are missing.

Business rule validation ensures that data adheres to application-specific rules. For example, meal logs cannot have future dates, duplicate meal logs within one hour are prevented, and exercise logs cannot have future dates or duplicates within two hours.

\subsection{SQL Injection Prevention}

SQL injection attacks are prevented by using Prisma ORM, which uses parameterized queries for all database operations. Parameterized queries separate SQL code from data, ensuring that user inputs are treated as data values rather than executable code. Prisma automatically escapes special characters in user inputs, preventing malicious SQL code from being executed.

\subsection{Cross-Site Scripting Prevention}

React provides built-in protection against XSS attacks by automatically escaping all values rendered in JSX. When displaying user-generated content, React converts special HTML characters to their entity equivalents, preventing malicious scripts from being executed. The application avoids using dangerouslySetInnerHTML, which bypasses React's XSS protection, except in cases where HTML content is explicitly trusted and sanitized.

\section{Frontend Design and User Experience}

The frontend is designed with user experience as a primary concern. The design follows several key principles to ensure usability and accessibility.

\subsection{Responsive Design}

The application uses a mobile-first responsive design approach. CSS styles are written first for mobile devices, then enhanced for larger screens using media queries. Tailwind CSS provides responsive utility classes that make it easy to apply different styles at different breakpoints. The application defines three main breakpoints: mobile for screens smaller than 768 pixels, tablet for screens between 768 and 1024 pixels, and desktop for screens larger than 1024 pixels.

Layout components adapt to screen size using flexible grid systems. On mobile devices, content is displayed in a single column to maximize readability. On tablets, content is arranged in two columns where appropriate. On desktop screens, content can span three or four columns, taking advantage of the available space.

Touch targets are sized appropriately for mobile devices. Buttons and interactive elements have minimum dimensions of 44 by 44 pixels, ensuring they can be easily tapped with a finger. Spacing between interactive elements prevents accidental taps on adjacent elements.

\subsection{Visual Feedback}

The application provides immediate visual feedback for all user actions. When users click buttons, the buttons change appearance to indicate they have been pressed. Loading states are indicated with spinner animations and disabled buttons, preventing users from submitting forms multiple times. Success and error messages appear as colored banners at the top of the screen, with green for success and red for errors.

Form inputs provide real-time validation feedback. As users type, the application validates their input and displays error messages below fields when validation fails. Valid inputs are indicated with green borders or checkmarks. This immediate feedback helps users correct errors before submitting forms.

\subsection{Accessibility}

The application follows web accessibility guidelines to ensure usability for all users. All interactive elements are keyboard accessible, allowing users to navigate the application without a mouse. Focus indicators clearly show which element currently has keyboard focus. Form labels are properly associated with their inputs using the htmlFor attribute, enabling screen readers to announce labels when inputs receive focus.

Color is not used as the only means of conveying information. For example, error messages include both red color and an error icon, ensuring that colorblind users can still identify errors. Sufficient color contrast is maintained between text and backgrounds, meeting WCAG AA standards for readability.

\section{Testing}

\subsection{Testing Approach}

The application was tested using a multi-level testing strategy covering unit tests, integration tests, and system tests. This comprehensive approach ensures that individual components work correctly in isolation, components interact properly when integrated, and the complete system meets all requirements.

\subsection{Test Cases and Results}

This section presents five major test cases that validate critical functionality of the NutriFit application. Each test case includes the test objective, input data, expected results, actual results, and pass/fail status.

\subsubsection{Test Case 1: User Registration and Authentication}

\textbf{Objective:} Verify that new users can register successfully and authenticate with correct credentials.

\textbf{Test Steps:}
\begin{enumerate}
    \item Navigate to registration page
    \item Enter name: "John Doe", email: "john@example.com", password: "SecurePass123"
    \item Submit registration form
    \item Verify JWT token is generated
    \item Logout and attempt login with same credentials
\end{enumerate}

\textbf{Expected Results:} User account created successfully, JWT token generated, login successful with correct credentials, login fails with incorrect credentials.

\textbf{Actual Results:} User registration completed successfully. Database record created with hashed password. JWT token generated with 7-day expiration. Login successful with correct credentials. Login rejected with error message for incorrect credentials.

\begin{figure}[H]
\centering
\includegraphics[width=0.85\textwidth]{test-case-1-registration.png}
\caption{Test Case 1: User Registration Test Result}
\label{fig:test1}
\end{figure}

Figure~\ref{fig:test1} shows the successful execution of the user registration test case. The screenshot demonstrates that the registration form accepts valid input, creates a user account in the database, and generates an authentication token.

\textbf{Status:} \textcolor{green}{PASSED} ✓

\subsubsection{Test Case 2: BMI Calculation and Health Profile}

\textbf{Objective:} Verify that BMI calculation produces correct results and health profiles are saved properly.

\textbf{Test Data:}
\begin{itemize}
    \item Weight: 75 kg, Height: 175 cm
    \item Age: 28, Gender: Male
    \item Activity Level: Moderate, Goal: Maintain Weight
\end{itemize}

\textbf{Expected Results:} BMI = 24.5 (Normal category), Daily calories ≈ 2400-2500 kcal.

\textbf{Actual Results:} BMI calculated as 24.5, correctly classified as "Normal Weight". Daily calorie requirement calculated as 2,456 kcal using Harris-Benedict equation. Health profile saved to database with all fields correctly stored.

\begin{figure}[H]
\centering
\includegraphics[width=0.85\textwidth]{test-case-2-bmi.png}
\caption{Test Case 2: BMI Calculation Test Result}
\label{fig:test2}
\end{figure}

Figure~\ref{fig:test2} displays the BMI calculation test results showing accurate computation of body mass index and proper categorization. The health profile form correctly stores all user data and performs real-time calculations.

\textbf{Status:} \textcolor{green}{PASSED} ✓

\subsubsection{Test Case 3: Diet Plan Generation with Allergen Safety}

\textbf{Objective:} Verify that diet plans exclude foods containing user allergens and meet calorie targets.

\textbf{Test Data:}
\begin{itemize}
    \item User allergens: Peanuts, Dairy
    \item Daily calorie target: 2000 kcal
    \item Health goal: Lose weight
\end{itemize}

\textbf{Expected Results:} Generated diet plan contains no foods with peanuts or dairy. Total daily calories approximately 2000 kcal distributed across 4 meals.

\textbf{Actual Results:} Diet plan generated successfully with breakfast (500 kcal), lunch (700 kcal), dinner (600 kcal), and snacks (200 kcal) totaling 2000 kcal. All foods verified to be free of peanuts and dairy allergens. Safety checker correctly filtered 12 unsafe foods from the catalog.

\begin{figure}[H]
\centering
\includegraphics[width=0.85\textwidth]{test-case-3-diet-plan.png}
\caption{Test Case 3: Diet Plan Generation with Allergen Safety Test Result}
\label{fig:test3}
\end{figure}

Figure~\ref{fig:test3} illustrates the diet plan generation test showing that the system successfully excludes allergen-containing foods and generates a balanced meal plan meeting the specified calorie target.

\textbf{Status:} \textcolor{green}{PASSED} ✓

\subsubsection{Test Case 4: Exercise Safety Validation}

\textbf{Objective:} Verify that high-impact exercises are flagged as unsafe for users with medical conditions.

\textbf{Test Data:}
\begin{itemize}
    \item User medical conditions: Heart Disease, Knee Problems
    \item Exercise: Running (High Impact)
\end{itemize}

\textbf{Expected Results:} Exercise flagged as unsafe with appropriate warning message.

\textbf{Actual Results:} Safety checker correctly identified running as high-impact exercise. Warning message displayed: "High-impact exercises not recommended for users with heart disease or joint problems. Consider low-impact alternatives like swimming or cycling." Workout plan generator excluded all high-impact exercises and selected only low-impact alternatives.

\begin{figure}[H]
\centering
\includegraphics[width=0.85\textwidth]{test-case-4-exercise-safety.png}
\caption{Test Case 4: Exercise Safety Validation Test Result}
\label{fig:test4}
\end{figure}

Figure~\ref{fig:test4} demonstrates the exercise safety validation feature showing warning messages for unsafe exercises based on user medical conditions. The system successfully prevents potentially harmful exercise recommendations.

\textbf{Status:} \textcolor{green}{PASSED} ✓

\subsubsection{Test Case 5: Goal Progress Tracking}

\textbf{Objective:} Verify that goal progress is calculated correctly and status indicators update appropriately.

\textbf{Test Data:}
\begin{itemize}
    \item Goal Type: Weight Loss
    \item Start Value: 85 kg, Target Value: 75 kg
    \item Current Value: 80 kg
    \item Deadline: 30 days from start
\end{itemize}

\textbf{Expected Results:} Progress = 50\%, Status = "On Track" (green indicator).

\textbf{Actual Results:} Progress calculated as 50\% ((85-80)/(85-75) × 100). Status indicator displayed in green with "On Track" label. Progress bar filled to 50\%. Remaining weight to lose: 5 kg. Days remaining: 15 days. System correctly updated progress when user logged new weight measurements.

\begin{figure}[H]
\centering
\includegraphics[width=0.85\textwidth]{test-case-5-goal-tracking.png}
\caption{Test Case 5: Goal Progress Tracking Test Result}
\label{fig:test5}
\end{figure}

Figure~\ref{fig:test5} shows the goal progress tracking test results with accurate percentage calculations and appropriate status indicators. The progress bar and statistics update correctly as users log their progress.

\textbf{Status:} \textcolor{green}{PASSED} ✓

\subsection{Test Summary}

All five test cases passed successfully, demonstrating that the application meets its functional requirements. The tests validated critical functionality including authentication, health calculations, safety mechanisms, and progress tracking. Performance testing showed that page load times are consistently under three seconds and API response times are under one second. Browser compatibility testing confirmed that the application works correctly in Chrome, Firefox, Safari, and Edge.


\section{Key Code Implementations}

This section presents essential code samples that demonstrate the core functionality of the NutriFit application. Each code sample is explained in detail to illustrate the implementation approach and design decisions.

\subsection{User Authentication}

The user registration process implements secure password hashing and JWT token generation as shown in Listing~\ref{code:registration}. The code first validates that all required fields are present in the request. It then checks the database to ensure the email address is not already registered, preventing duplicate accounts. The password is hashed using bcrypt with a salt factor of 10, which provides strong protection against brute force attacks. After creating the user record in the database, the system generates a JWT token containing the user's ID and email, which is returned to the client for subsequent authenticated requests.

\begin{lstlisting}[language=JavaScript, caption=User Registration Implementation, label=code:registration]
// app/api/auth/register/route.ts
export async function POST(request: Request) {
  const { name, email, password } = await request.json();
  
  // Validate required fields
  if (!name || !email || !password) {
    return NextResponse.json(
      { error: 'All fields required' }, 
      { status: 400 }
    );
  }
  
  // Check for existing user
  const existingUser = await prisma.user.findUnique({
    where: { email }
  });
  
  if (existingUser) {
    return NextResponse.json(
      { error: 'Email already registered' }, 
      { status: 400 }
    );
  }
  
  // Hash password and create user
  const passwordHash = await bcrypt.hash(password, 10);
  const user = await prisma.user.create({
    data: { name, email, passwordHash }
  });
  
  // Generate JWT token
  const token = await generateToken({
    userId: user.id,
    email: user.email
  });
  
  return NextResponse.json({ token, user });
}
\end{lstlisting}

The registration code shown in Listing~\ref{code:registration} demonstrates several important security practices. Input validation prevents empty fields from being processed. The unique email check prevents account duplication. Password hashing ensures that plain text passwords are never stored in the database. The JWT token provides stateless authentication without requiring server-side session storage.

\subsection{Health Calculations}

The BMI and calorie calculation functions implement scientifically validated formulas as presented in Listing~\ref{code:calculations}. The calculateBMI function converts height from centimeters to meters, then applies the standard BMI formula of weight divided by height squared. The result is classified into one of four categories based on World Health Organization guidelines. The calculateDailyCalories function implements the Harris-Benedict equation, which calculates Basal Metabolic Rate differently for males and females based on their weight, height, and age. The BMR is then multiplied by an activity level factor to determine Total Daily Energy Expenditure. Finally, the calorie value is adjusted based on the user's health goal, subtracting 500 calories for weight loss or adding 500 calories for weight gain.

\begin{lstlisting}[language=JavaScript, caption=BMI and Calorie Calculations, label=code:calculations]
// lib/safetyCheck.ts

// BMI Calculation
export function calculateBMI(weight, height) {
  if (!weight || !height) {
    return { bmi: null, error: 'Missing data' };
  }
  
  const heightInMeters = height / 100;
  const bmi = weight / (heightInMeters * heightInMeters);
  
  let category;
  if (bmi < 18.5) category = 'Underweight';
  else if (bmi < 25) category = 'Normal';
  else if (bmi < 30) category = 'Overweight';
  else category = 'Obese';
  
  return { bmi: Math.round(bmi * 10) / 10, category };
}

// Daily Calorie Calculation (Harris-Benedict)
export function calculateDailyCalories(
  weight, height, age, gender, activityLevel, healthGoals
) {
  let bmr;
  if (gender === 'male') {
    bmr = 88.362 + (13.397 * weight) + 
          (4.799 * height) - (5.677 * age);
  } else {
    bmr = 447.593 + (9.247 * weight) + 
          (3.098 * height) - (4.330 * age);
  }
  
  const multipliers = {
    sedentary: 1.2, light: 1.375, moderate: 1.55,
    active: 1.725, very_active: 1.9
  };
  
  let tdee = bmr * multipliers[activityLevel];
  
  if (healthGoals === 'lose_weight') tdee -= 500;
  else if (healthGoals === 'gain_weight') tdee += 500;
  
  return Math.round(tdee);
}
\end{lstlisting}

The calculation functions in Listing~\ref{code:calculations} are critical for providing accurate health recommendations. The BMI calculation provides users with an objective measure of their body composition. The calorie calculation ensures that diet plans are tailored to individual metabolic needs and health goals. Both functions include error handling to manage invalid inputs gracefully.

\subsection{Safety Validation}

The exercise safety checker implements medical condition validation as demonstrated in Listing~\ref{code:safety}. The function maintains arrays of high-risk medical conditions and joint-related issues. It checks whether the user has any high-risk conditions such as heart disease or hypertension combined with high-impact exercises. It also checks for joint problems combined with high-impact activities. When unsafe combinations are detected, the function returns a safety status of false along with an appropriate warning message that explains the risk and suggests alternatives.

\begin{lstlisting}[language=JavaScript, caption=Exercise Safety Validation, label=code:safety]
// lib/safetyCheck.ts
export function checkExerciseSafety(impactLevel, diseases) {
  const highRisk = ['Heart Disease', 'Hypertension'];
  const jointIssues = ['Arthritis', 'Knee Problems', 
                       'Back Pain'];
  
  if (diseases.some(d => highRisk.includes(d)) && 
      impactLevel === 'High') {
    return { 
      safe: false, 
      warning: 'High-intensity not recommended' 
    };
  }
  
  if (diseases.some(d => jointIssues.includes(d)) && 
      impactLevel === 'High') {
    return { 
      safe: false, 
      warning: 'High-impact may worsen joint pain' 
    };
  }
  
  return { safe: true };
}
\end{lstlisting}

The safety validation code in Listing~\ref{code:safety} demonstrates the application's commitment to user safety. By checking medical conditions against exercise characteristics, the system prevents potentially harmful recommendations. This proactive approach protects users from injury and ensures that the application provides responsible health guidance.

\subsection{Diet Plan Generation}

The diet plan generator orchestrates multiple steps to create personalized meal plans as shown in Listing~\ref{code:dietplan}. The function first calculates the user's daily calorie needs based on their health profile. It then distributes these calories across four meals following nutritional guidelines: breakfast receives 25 percent, lunch 35 percent, dinner 30 percent, and snacks 10 percent. The function retrieves all available foods from the database and filters out any foods containing allergens specified in the user's profile. For each meal, it calls a helper function to generate balanced food combinations. Finally, the complete diet plan is saved to the database and returned to the user.

\begin{lstlisting}[language=JavaScript, caption=Diet Plan Generator, label=code:dietplan]
// lib/planGenerator.ts
export async function generateDietPlan(userId, profile) {
  // Calculate daily calories
  const dailyCalories = calculateDailyCalories(
    profile.weight, profile.height, profile.age,
    profile.gender, profile.activityLevel, 
    profile.healthGoals
  );
  
  // Distribute calories across meals
  const distribution = {
    breakfast: dailyCalories * 0.25,
    lunch: dailyCalories * 0.35,
    dinner: dailyCalories * 0.30,
    snacks: dailyCalories * 0.10
  };
  
  // Get safe foods (exclude allergens)
  const allergens = JSON.parse(profile.allergens || '[]');
  const allFoods = await prisma.foodItem.findMany();
  const safeFoods = allFoods.filter(food => {
    const foodAllergens = JSON.parse(food.allergens || '[]');
    return !foodAllergens.some(a => allergens.includes(a));
  });
  
  // Generate balanced meals
  const meals = {
    breakfast: generateMeal(safeFoods, 'breakfast'),
    lunch: generateMeal(safeFoods, 'lunch'),
    dinner: generateMeal(safeFoods, 'dinner'),
    snacks: generateMeal(safeFoods, 'snacks')
  };
  
  // Save to database
  return await prisma.dietPlan.create({
    data: { 
      userId, 
      meals: JSON.stringify(meals), 
      calories: dailyCalories 
    }
  });
}
\end{lstlisting}

The diet plan generation code in Listing~\ref{code:dietplan} illustrates the integration of multiple system components. The function combines health calculations, database queries, safety filtering, and meal composition algorithms to produce personalized nutrition plans. The calorie distribution follows established nutritional guidelines to ensure balanced energy intake throughout the day.

\section{Database Schema}

The database schema for NutriFit consists of eleven interconnected tables designed following normalization principles. This section presents the complete Prisma schema definition and explains the relationships between entities.

\subsection{Schema Overview}

The complete database schema is presented in Listing~\ref{code:schema}. The schema defines eleven models representing different entities in the system. The User model serves as the central entity with one-to-one relationship to HealthProfile and one-to-many relationships to all other user-specific entities. Foreign key constraints with CASCADE delete behavior ensure that when a user is deleted, all related records are automatically removed. Unique constraints prevent duplicate data such as duplicate email addresses or multiple water intake records for the same user on the same day. Indexes on frequently queried columns improve query performance for operations like retrieving logs by user and date.

\begin{lstlisting}[language=SQL, caption=Complete Prisma Database Schema, label=code:schema]
// prisma/schema.prisma

generator client {
  provider = "prisma-client-js"
}

datasource db {
  provider = "sqlite"
  url      = env("DATABASE_URL")
}

model User {
  id            Int             @id @default(autoincrement())
  name          String
  email         String          @unique
  passwordHash  String
  createdAt     DateTime        @default(now())
  
  healthProfile HealthProfile?
  mealLogs      MealLog[]
  exerciseLogs  ExerciseLog[]
  dietPlans     DietPlan[]
  workoutPlans  WorkoutPlan[]
  goals         Goal[]
  waterIntake   WaterIntake[]
  favorites     Favorite[]
}

model HealthProfile {
  id            Int      @id @default(autoincrement())
  userId        Int      @unique
  age           Int?
  weight        Float?
  height        Float?
  gender        String?  @default("male")
  activityLevel String?  @default("moderate")
  healthGoals   String?  @default("maintain")
  diseases      String?
  allergens     String?
  createdAt     DateTime @default(now())
  updatedAt     DateTime @updatedAt
  
  user          User     @relation(fields: [userId], 
                         references: [id], onDelete: Cascade)
}

model Goal {
  id            Int      @id @default(autoincrement())
  userId        Int
  goalType      String
  targetValue   Float
  currentValue  Float?
  deadline      DateTime?
  status        String   @default("active")
  createdAt     DateTime @default(now())
  
  user          User     @relation(fields: [userId], 
                         references: [id], onDelete: Cascade)
  
  @@index([userId, status])
}

model DietPlan {
  id        Int      @id @default(autoincrement())
  userId    Int
  meals     String
  calories  Int?
  createdAt DateTime @default(now())
  
  user      User     @relation(fields: [userId], 
                     references: [id], onDelete: Cascade)
}

model MealLog {
  id        Int      @id @default(autoincrement())
  userId    Int
  date      DateTime @default(now())
  meals     String
  createdAt DateTime @default(now())
  
  user      User     @relation(fields: [userId], 
                     references: [id], onDelete: Cascade)
  
  @@index([userId, date])
}

model WaterIntake {
  id        Int      @id @default(autoincrement())
  userId    Int
  date      DateTime @default(now())
  glasses   Int      @default(0)
  goal      Int      @default(8)
  createdAt DateTime @default(now())
  
  user      User     @relation(fields: [userId], 
                     references: [id], onDelete: Cascade)
  
  @@unique([userId, date])
}

model FoodItem {
  id        Int      @id @default(autoincrement())
  name      String
  calories  Int
  allergens String?
  createdAt DateTime @default(now())
}

model ExerciseItem {
  id          Int      @id @default(autoincrement())
  name        String
  impactLevel String
  createdAt   DateTime @default(now())
}
\end{lstlisting}

The database schema shown in Listing~\ref{code:schema} demonstrates careful design for data integrity and performance. The User table stores authentication credentials with hashed passwords. The HealthProfile table maintains a one-to-one relationship with User, ensuring each user has at most one health profile. The Goal, DietPlan, MealLog, and WaterIntake tables have one-to-many relationships with User, allowing users to have multiple records of each type. The FoodItem and ExerciseItem tables serve as catalog tables that are referenced by other entities but do not have direct foreign key relationships. The schema uses appropriate data types for each field, with integers for IDs and counts, floats for measurements like weight and height, strings for text data, and DateTime for timestamps. Default values ensure that fields have sensible initial values when records are created.
